\documentclass{zhvt-classic}

% 使用本地源流明体（GenRyuMin2 TW）按文件路径加载，不依赖系统字体名
% 版框（内框／界栏／书口侧栏／鱼尾）改用与版心印章同一印泥棕色：
%   印章图片 yshs.png 实色 = RGB(135,68,52) = HTML 874434
\usepackage{xcolor}
\usepackage{tikz}
% sealbrown 双注册：xcolor 供 TikZ/pgf，L3 供 cls 版框
\definecolor{sealbrown}{HTML}{874434}
\ExplSyntaxOn
\color_set:nnn { sealbrown } { HTML } { 874434 }
\ExplSyntaxOff

\zhvtset{
  grid color       = sealbrown,
  main font        = GenRyuMin2TW-SB,
  main font path   = C:/Users/Administrator/texmf/fonts/opentype/public/GenRyuMin2TW/,
jiazhu font      = GenRyuMin2TW-M,
  jiazhu font path = C:/Users/Administrator/texmf/fonts/opentype/public/GenRyuMin2TW/,
  % 页码框容纳 4 字：正常页「九十七」+前置空 1 格=4 行恰好；
  % 100+ 页「一百零一」共 5 行超出，末字「一/二/三」被框裁掉（2026-10-04 实测），故放宽到 5
  shukou page number chars = 5,
  % 书口书名框高从默认 6 字减到 4 字，把省出的 2 字高拨给 chapter 框，
  % 使 6 字书名+卷名放进 9 字高的 chapter 框后不再压线（2026-10-05）
  shukou title chars = 4,
}

\title{史记}[吴王濞列传]
\author{汉·司马迁}
\maker{闲叟漫记·重製}


% ===== 索隐述赞样式（2026-11-08 四次定稿 · 愚大人亲手改定参数）=====
% ⚠ 竖排前提：自定义字体家族必须带 RawFeature={vertical:+vert}；
%   simhei.ttf 加 +vhal 会整体加载失败静默回退，只写 +vert；simfang 同理。
\newCJKfontfamily{\shuzanfont}[Path=C:/Windows/Fonts/,RawFeature={vertical:+vert}]{simhei.ttf}
\newCJKfontfamily{\shuzanbody}[Path=C:/Windows/Fonts/,RawFeature={vertical:+vert}]{simfang.ttf}
% 题头：152x30pt 圆角 6pt 棕底白字；baseline (0,2) 使整框落入界行格子不压版框线；
% 白字节点 x 不补偿、y +8.2pt（像素级墨迹实测 dx/dy≈0，L/R 留白约 20pt）。
\newcommand{\shuzanhead}[1]{%
  \begin{tikzpicture}[baseline={(0,2)}]
    \fill[sealbrown, rounded corners=6pt]
      (-76pt,-15pt) rectangle (76pt,15pt);
    \node[text=white, inner sep=0pt, anchor=center] at (0pt,8.2pt)
      {\shuzanfont\fontsize{23}{23}\selectfont #1};
  \end{tikzpicture}%
}
% ★ 修正书口卷名位置：chapter 框起始位置改回 cls 默认值，避免全局影响。
%   只在每个卷首页通过 \juan 命令里的空命令微调该页位置（安全，不会跨页污染）。
\ExplSyntaxOn
\AtBeginDocument{
  \dim_set:Nn \l__zhvt_shukou_chapter_title_start_dim
    { \l__zhvt_shukou_title_height_dim + \l__zhvt_font_size_dim * 2 }
  \dim_set:Nn \l__zhvt_shukou_page_number_start_dim
    { \l__zhvt_shukou_chapter_title_start_dim + \l__zhvt_shukou_chapter_title_height_dim }
}
\ExplSyntaxOff

\begin{document}

% ★★ 版心（书口）排布覆写：把书名放到鱼尾之下（2026-10-04）
%
%   cls 默认版心自上而下：书名（title 静态框，在鱼尾之上）→ 鱼尾 → 篇名（chapter 动态框）→ 页码
%   本模板改为：          鱼尾（位置不动）→ 书名 → 篇名 → 页码分割线 → 页码
%
%   手法：把 cls 放在 title 静态框里的书名，改放进 chapter 动态框（鱼尾之下）；
%         title 静态框留空。书名取 \l__zhvt_title_str，即 \title{...} 的第一参数，
%         所以换书时只改 \title，不必在此重复写书名。
%
%   ⚠ 三条必须遵守的写法（2026-10-04 实测）：
%     ① 汉字必须用 \hbox{...} 包住。flowfram 会把框内容装进一个宽 14.4mm（= 一个 line_sep
%        = 书口栏宽）的纵向 minipage 并按该宽度折行，而整页靠 PDF /Rotate 90 显示、
%        框内文字本身不变向 —— 不包 \hbox 时「围炉夜话」4 字会被逐字折成 4 行，
%        在版面上就成了横排 4 字并横向溢出版框。
%     ② 竖排下**一个 \hbox 就是一整列**。书名与篇名要同处一列，必须放进同一个 \hbox；
%        分成两个 \hbox 时，后一个会被当成新的一列、甩到隔壁栏去
%        （实测「序」跑到书口栏外侧 14.4mm 的正文列上）。中间的 \quad 在框内是横向空白，
%        显示出来正好是书名与篇名之间的纵向间隔。
%     ③ ★ 内容盒要写成 **与框等宽的定宽盒** `\hbox to \l__zhvt_shukou_chapter_title_height_dim{…\hfil}`。
%        实测：chapter 框是把内容盒沿纵向**居中**摆放的，内容宽度一变，书名就跟着漂：
%        自然宽度的 \hbox（书名 38.4mm）落在 x=157.2mm，加篇名后（57.6mm）落在 x=152.4mm，
%        目录页（无篇名）与正文页的书名就差半格。做成等宽盒后，无论有没有篇名，
%        内容固定从框顶起排，实测两页均落在 x=142.8mm，不再漂。
\ExplSyntaxOn
% 把书口 chapter 框（书名+篇名）起始位置下移 1 字高，避免 4 字卷名顶到鱼尾
% ⚠ cls 的 \__zhvt_calculate_dimensions: 在 \AtBeginDocument 里执行，会重算该 dim；
%   必须把覆写挂到 begin document 之后，否则会被二次覆盖（2026-10-05 实测）。
% (旧方案已移到导言区，直接修改 chapter 框起始位置)
% 卷标题命令 #1=书口篇名(鱼尾) #2=正文卷名（含卷序，卷名内空格用 ~）
%
% ★★ 目录页两项修订（2026-10-05，愚大人指出）★★
%   ① 页码起点统一对齐：
%      原状 —— 目录项写成 \par\hspace{4em}卷名\hspace{4em}页码，页码起点 = 卷名实宽 + 4em。
%      卷名实宽随字数变：「卷一~梁惠王上」6 字 +~、「卷七~离娄上」5 字 +~，
%      差 1 个字位 → 卷七~卷十 4 条的页码整体比其余 10 条高 27.21pt（一个字），看着上下不齐。
%      改法 —— 卷名补全角空格（宽 = 1 字 = 27.215pt）补足到 **7 字符位**再写进目录，
%      14 条页码起点遂全部落在同一条竖线上（起点对齐；末字随页码字数自然参差）。
%      ⚠ 只对目录里的副本 \l_tmpc_tl 补，正文卷名仍用 \l_tmpa_tl 原样 —— 正文像素不受影响。
%   ② 卷名整体上移半字位：前导 \hspace 由 4em 改 3.5em。
%      实测原状目录块在版框内偏下（天头留白 139.6pt / 地脚留白 110.7pt），
%      上移 0.5em = 13.6pt 后上下基本均分，视觉居中。
\NewDocumentCommand { \juan } { mm } {
	\clearpage
	\tl_set:Nn \l__zhvt_chapter_title_tl { #1 }
	\tl_set:Nn \l_tmpa_tl { #2 }
	\tl_set_eq:NN \l_tmpc_tl \l_tmpa_tl
	\int_step_inline:nn { 7 - \str_count:n { #2 } }
	  { \tl_put_right:Nn \l_tmpc_tl { 　 } }
	\tl_set:Nx \l_tmpb_tl { \zhnum { page } }
	\tl_put_left:Nn \l_tmpb_tl { \hspace { 4em } }
	\tl_put_left:Nn \l_tmpb_tl { \l_tmpc_tl }
	\tl_put_left:Nn \l_tmpb_tl { \par \hspace { 3.5em } }
	\addtocontents { toc } { \l_tmpb_tl }
	\skip_horizontal:n { \l__zhvt_font_size_dim * 4 }
	\l_tmpa_tl
	\par
	\thispagestyle { empty }
}
\cs_set_protected:Npn \__zhvt_set_shukou_content: {
	% ★ 鱼尾上方留白（title 静态框）嵌入闲叟钤印：
	%   整页 /Rotate 90，故此处把图用 \rotatebox{90} 转正为竖章；
	%   尺寸卡在版心条幅内（书内约 7mm 宽 × 33.8mm 高，留余量不压鱼尾）。
	\setstaticcontents * { title } {
		\vbox_to_ht:nn { \l__zhvt_baseline_skip_dim } {
			\vskip 14.2pt
			\hbox_to_wd:nn { \l__zhvt_shukou_title_height_dim } {
				\hfil
				\rotatebox [ origin = c ] { 90 } { \includegraphics [ width = 7mm , height = 33.8mm ] { yshs.png } }
				\hfil
			}
		}
	}
% ★ 版心（书口）统一改法（2026-11-16 愚大人定）：18pt＋框内居中＋书口文字同版框色
%   ① 原 27.2pt 时 7 字长篇名（如「廉颇蔺相如列传」）末字溢出 chapter 框、压住页码；
%      改 \fontsize{18pt}{18pt}，前后 \hfil 令字串在框内居中（不随篇名字数漂移）。
%   ② \vbox{\vskip 8.2pt} 定量下移，抵消缩字号后栏宽方向 8.2pt 上偏（像素实测 y 中心 543.3 = 原版）。
%      ⚠ \raisebox / \vspace* 在 flowfram 框内无效；换字号须同步换补偿：16pt→10.1、19pt→7.4。
%   ③ chapter 与 page~number 均置 \color{sealbrown}（HTML 874434），与版框／界行同色；正文仍黑。
%   ⚠ 页码字位、钤印、鱼尾、界行、正文一律未动；五部合订本同此口径。
	\setdynamiccontents * { chapter }
	{ \color{sealbrown} \vbox { \vskip 8.2pt \hbox to \l__zhvt_shukou_chapter_title_height_dim
		{ \hfil \fontsize{18pt}{18pt}\selectfont \l__zhvt_title_str \quad \l__zhvt_chapter_title_tl \hfil } } }
	\setdynamiccontents * { page~number } { \color{sealbrown} \quad \zhnum { page } }
}
\__zhvt_set_shukou_content:
\ExplSyntaxOff

\maketitle

% 单篇，无多卷目录，不渲染目录页

\juan{吴王濞列传}{卷一百零六~吴王濞列传第四十六}
\hspace*{2em}吴王濞者。高帝兄刘仲之子也。高帝已定天下七年。立刘仲为代王。而匈奴攻代。刘仲不能坚守。弃国亡。间行走雒阳。自归天子。天子为骨肉故。不忍致法。废以为郃阳侯。高帝十一年秋。淮南王英布反。东并荆地。劫其国兵。西度淮。击楚。高帝自将往诛之。刘仲子沛侯濞年二十。有气力。以骑将从破布军蕲西。会甀。布走。荆王刘贾为布所杀。无后。上患吴。会稽轻悍。无壮王以填之。诸子少。乃立濞於沛为吴王。王三郡五十三城。已拜受印。高帝召濞相之。谓曰。﹁若状有反相。﹂心独悔。业已拜。因拊其背。告曰。﹁汉后五十年东南有乱者。岂若邪。然天下同姓为一家也。慎无反。﹂濞顿首曰。﹁不敢。﹂\\
\hspace*{2em}会孝惠。高后时。天下初定。郡国诸侯各务自拊循其民。吴有豫章郡铜山。濞则招致天下亡命者铸钱。煮海水为盐。以故无赋。国用富饶。\\
\hspace*{2em}孝文时。吴太子入见。得侍皇太子饮博。吴太子师傅皆楚人。轻悍。又素骄。博。争道。不恭。皇太子引博局提吴太子。杀之。於是遣其丧归葬。至吴。吴王愠曰。﹁天下同宗。死长安即葬长安。何必来葬为。﹂复遣丧之长安葬。吴王由此稍失籓臣之礼。称病不朝。京师知其以子故称病不朝。验问实不病。诸吴使来。辄系责治之。吴王恐。为谋滋甚。及后使人为秋请。上复责问吴使者。使者对曰。﹁王实不病。汉系治使者数辈。以故遂称病。且夫﹃察见渊中鱼。不祥﹄。今王始诈病。及觉。见责急。愈益闭。恐上诛之。计乃无聊。唯上弃之而与更始。﹂於是天子乃赦吴使者归之。而赐吴王几杖。老。不朝。吴得释其罪。谋亦益解。然其居国以铜盐故。百姓无赋。卒践更。辄与平贾。岁时存问茂材。赏赐闾里。佗郡国吏欲来捕亡人者。讼共禁弗予。如此者四十馀年。以故能使其众。\\
\hspace*{2em}晁错为太子家令。得幸太子。数从容言吴过可削。数上书说孝文帝。文帝宽。不忍罚。以此吴日益横。及孝景帝即位。错为御史大夫。说上曰。﹁昔高帝初定天下。昆弟少。诸子弱。大封同姓。故王孽子悼惠王王齐七十馀城。庶弟元王王楚四十馀城。兄子濞王吴五十馀城。封三庶孽。分天下半。今吴王前有太子之郄。诈称病不朝。於古法当诛。文帝弗忍。因赐几杖。德至厚。当改过自新。乃益骄溢。即山铸钱。煮海水为盐。诱天下亡人。谋作乱。今削之亦反。不削之亦反。削之。其反亟。祸小。不削。反迟。祸大。﹂三年冬。楚王朝。晁错因言楚王戊往年为薄太后服。私奸服舍。请诛之。诏赦。罚削东海郡。因削吴之豫章郡。会稽郡。及前二年赵王有罪。削其河间郡。胶西王卬以卖爵有奸。削其六县。\\
\hspace*{2em}汉廷臣方议削吴。吴王濞恐削地无已。因以此发谋。欲举事。念诸侯无足与计谋者。闻胶西王勇。好气。喜兵。诸齐皆惮畏。於是乃使中大夫应高誂胶西王。无文书。口报曰。﹁吴王不肖。有宿夕之忧。不敢自外。使喻其驩心。﹂王曰。﹁何以教之。﹂高曰。﹁今者主上兴於奸。饰於邪臣。好小善。听谗贼。擅变更律令。侵夺诸侯之地。徵求滋多。诛罚良善。日以益甚。里语有之。﹃舐糠及米﹄。吴与胶西。知名诸侯也。一时见察。恐不得安肆矣。吴王身有内病。不能朝请二十馀年。尝患见疑。无以自白。今胁肩累足。犹惧不见释。窃闻大王以爵事有適。所闻诸侯削地。罪不至此。此恐不得削地而已。﹂王曰。﹁然。有之。子将柰何。﹂高曰。﹁同恶相助。同好相留。同情相成。同欲相趋。同利相死。今吴王自以为与大王同忧。原因时循理。弃躯以除患害於天下。亿亦可乎。﹂王瞿然骇曰。﹁寡人何敢如是。今主上虽急。固有死耳。安得不戴。﹂高曰。﹁御史大夫晁错。荧惑天子。侵夺诸侯。蔽忠塞贤。朝廷疾怨。诸侯皆有倍畔之意。人事极矣。彗星出。蝗蟲数起。此万世一时。而愁劳圣人之所以起也。故吴王欲内以晁错为讨。外随大王后车。彷徉天下。所乡者降。所指者下。天下莫敢不服。大王诚幸而许之一言。则吴王率楚王略函谷关。守荥阳敖仓之粟。距汉兵。治次舍。须大王。大王有幸而临之。则天下可并。两主分割。不亦可乎。﹂王曰。﹁善。﹂高归报吴王。吴王犹恐其不与。乃身自为使。使於胶西。面结之。\\
\hspace*{2em}胶西群臣或闻王谋。谏曰。﹁承一帝。至乐也。今大王与吴西乡。弟令事成。两主分争。患乃始结。诸侯之地不足为汉郡什二。而为畔逆以忧太后。非长策也。﹂王弗听。遂发使约齐。菑川。胶东。济南。济北。皆许诺。而曰﹁城阳景王有义。攻诸吕。勿与。事定分之耳﹂。\\
\hspace*{2em}诸侯既新削罚。振恐。多怨晁错。及削吴会稽。豫章郡书至。则吴王先起兵。胶西正月丙午诛汉吏二千石以下。胶东。菑川。济南。楚。赵亦然。遂发兵西。齐王后悔。饮药自杀。畔约。济北王城坏未完。其郎中令劫守其王。不得发兵。胶西为渠率。胶东。菑川。济南共攻围临菑。赵王遂亦反。阴使匈奴与连兵。\\
\hspace*{2em}七国之发也。吴王悉其士卒。下令国中曰。﹁寡人年六十二。身自将。少子年十四。亦为士卒先。诸年上与寡人比。下与少子等者。皆发。﹂发二十馀万人。南使闽越。东越。东越亦发兵从。\\
\hspace*{2em}孝景帝三年正月甲子。初起兵於广陵。西涉淮。因并楚兵。发使遗诸侯书曰。﹁吴王刘濞敬问胶西王。胶东王。菑川王。济南王。赵王。楚王。淮南王。衡山王。庐江王。故长沙王子。幸教寡人。以汉有贼臣。无功天下。侵夺诸侯地。使吏劾系讯治。以僇辱之为故。不以诸侯人君礼遇刘氏骨肉。绝先帝功臣。进任奸宄。诖乱天下。欲危社稷。陛下多病志失。不能省察。欲举兵诛之。谨闻教。敝国虽狭。地方三千里。人虽少。精兵可具五十万。寡人素事南越三十馀年。其王君皆不辞分其卒以随寡人。又可得三十馀万。寡人虽不肖。原以身从诸王。越直长沙者。因王子定长沙以北。西走蜀。汉中。告越。楚王。淮南三王。与寡人西面。齐诸王与赵王定河间。河内。或入临晋关。或与寡人会雒阳。燕王。赵王固与胡王有约。燕王北定代。云中。抟胡众入萧关。走长安。匡正天子。以安高庙。原王勉之。楚元王子。淮南三王或不沐洗十馀年。怨入骨髓。欲一有所出之久矣。寡人未得诸王之意。未敢听。今诸王苟能存亡继绝。振弱伐暴。以安刘氏。社稷之所原也。敝国虽贫。寡人节衣食之用。积金钱。脩兵革。聚穀食。夜以继日。三十馀年矣。凡为此。原诸王勉用之。能斩捕大将者。赐金五千斤。封万户。列将。三千斤。封五千户。裨将。二千斤。封二千户。二千石。千斤。封千户。千石。五百斤。封五百户。皆为列侯。其以军若城邑降者。卒万人。邑万户。如得大将。人户五千。如得列将。人户三千。如得裨将。人户千。如得二千石。其小吏皆以差次受爵金。佗封赐皆倍军法。其有故爵邑者。更益勿因。原诸王明以令士大夫。弗敢欺也。寡人金钱在天下者往往而有。非必取於吴。诸王日夜用之弗能尽。有当赐者告寡人。寡人且往遗之。敬以闻。﹂\\
\hspace*{2em}七国反书闻天子。天子乃遣太尉条侯周亚夫将三十六将军。往击吴楚。遣曲周侯郦寄击赵。将军栾布击齐。大将军窦婴屯荥阳。监齐赵兵。\\
\hspace*{2em}吴楚反书闻。兵未发。窦婴未行。言故吴相袁盎。盎时家居。诏召入见。上方与晁错调兵笇军食。上问袁盎曰。﹁君尝为吴相。知吴臣田禄伯为人乎。今吴楚反。於公何如。﹂对曰。﹁不足忧也。今破矣。﹂上曰。﹁吴王即山铸钱。煮海水为盐。诱天下豪桀。白头举事。若此。其计不百全。岂发乎。何以言其无能为也。﹂袁盎对曰。﹁吴有铜盐利则有之。安得豪桀而诱之。诚令吴得豪桀。亦且辅王为义。不反矣。吴所诱皆无赖子弟。亡命铸钱奸人。故相率以反。﹂晁错曰。﹁袁盎策之善。﹂上问曰。﹁计安出。﹂盎对曰。﹁原屏左右。﹂上屏人。独错在。盎曰。﹁臣所言。人臣不得知也。﹂乃屏错。错趋避东厢。恨甚。上卒问盎。盎对曰。﹁吴楚相遗书。曰﹃高帝王子弟各有分地。今贼臣晁错擅適过诸侯。削夺之地﹄。故以反为名。西共诛晁错。复故地而罢。方今计独斩晁错。发使赦吴楚七国。复其故削地。则兵可无血刃而俱罢。﹂於是上嘿然良久。曰。﹁顾诚何如。吾不爱一人以谢天下。﹂盎曰。﹁臣愚计无出此。原上孰计之。﹂乃拜盎为太常。吴王弟子德侯为宗正。盎装治行。后十馀日。上使中尉召错。绐载行东市。错衣朝衣斩东市。则遣袁盎奉宗庙。宗正辅亲戚。使告吴如盎策。至吴。吴楚兵已攻梁壁矣。宗正以亲故。先入见。谕吴王使拜受诏。吴王闻袁盎来。亦知其欲说己。笑而应曰。﹁我已为东帝。尚何谁拜。﹂不肯见盎而留之军中。欲劫使将。盎不肯。使人围守。且杀之。盎得夜出。步亡去。走梁军。遂归报。\\
\hspace*{2em}条侯将乘六乘传。会兵荥阳。至雒阳。见剧孟。喜曰。﹁七国反。吾乘传至此。不自意全。又以为诸侯已得剧孟。剧孟今无动。吾据荥阳。以东无足忧者。﹂至淮阳。问父绛侯故客邓都尉曰。﹁策安出。﹂客曰。﹁吴兵锐甚。难与争锋。楚兵轻。不能久。方今为将军计。莫若引兵东北壁昌邑。以梁委吴。吴必尽锐攻之。将军深沟高垒。使轻兵绝淮泗口。塞吴饟道。彼吴梁相敝而粮食竭。乃以全彊制其罢极。破吴必矣。﹂条侯曰。﹁善。﹂从其策。遂坚壁昌邑南。轻兵绝吴饟道。\\
\hspace*{2em}吴王之初发也。吴臣田禄伯为大将军。田禄伯曰。﹁兵屯聚而西。无佗奇道。难以就功。臣原得五万人。别循江淮而上。收淮南。长沙。入武关。与大王会。此亦一奇也。﹂吴王太子谏曰。﹁王以反为名。此兵难以藉人。藉人亦且反王。柰何。且擅兵而别。多佗利害。未可知也。徒自损耳。﹂吴王即不许田禄伯。\\
\hspace*{2em}吴少将桓将军说王曰。﹁吴多步兵。步兵利险。汉多车骑。车骑利平地。原大王所过城邑不下。直弃去。疾西据雒阳武库。食敖仓粟。阻山河之险以令诸侯。虽毋入关。天下固已定矣。即大王徐行。留下城邑。汉军车骑至。驰入梁楚之郊。事败矣。﹂吴王问诸老将。老将曰。﹁此少年推锋之计可耳。安知大虑乎。﹂於是王不用桓将军计。\\
\hspace*{2em}吴王专并将其兵。未度淮。诸宾客皆得为将。校尉。候。司马。独周丘不得用。周丘者。下邳人。亡命吴。酤酒无行。吴王濞薄之。弗任。周丘上谒。说王曰。﹁臣以无能。不得待罪行间。臣非敢求有所将。原得王一汉节。必有以报王。﹂王乃予之。周丘得节。夜驰入下邳。下邳时闻吴反。皆城守。至传舍。召令。令入户。使从者以罪斩令。遂召昆弟所善豪吏告曰。﹁吴反兵且至。至。屠下邳不过食顷。今先下。家室必完。能者封侯矣。﹂出乃相告。下邳皆下。周丘一夜得三万人。使人报吴王。遂将其兵北略城邑。比至城阳。兵十馀万。破城阳中尉军。闻吴王败走。自度无与共成功。即引兵归下邳。未至。疽发背死。\\
\hspace*{2em}二月中。吴王兵既破。败走。於是天子制诏将军曰。﹁盖闻为善者。天报之以福。为非者。天报之以殃。高皇帝亲表功德。建立诸侯。幽王。悼惠王绝无后。孝文皇帝哀怜加惠。王幽王子遂。悼惠王子卬等。令奉其先王宗庙。为汉籓国。德配天地。明并日月。吴王濞倍德反义。诱受天下亡命罪人。乱天下币。称病不朝二十馀年。有司数请濞罪。孝文皇帝宽之。欲其改行为善。今乃与楚王戊。赵王遂。胶西王卬。济南王辟光。菑川王贤。胶东王雄渠约从反。为逆无道。起兵以危宗庙。贼杀大臣及汉使者。迫劫万民。夭杀无罪。烧残民家。掘其丘冢。甚为暴虐。今卬等又重逆无道。烧宗庙。卤御物。朕甚痛之。朕素服避正殿。将军其劝士大夫击反虏。击反虏者。深入多杀为功。斩首捕虏比三百石以上者皆杀之。无有所置。敢有议诏及不如诏者。皆要斩。﹂\\
\hspace*{2em}初。吴王之度淮。与楚王遂西败棘壁。乘胜前。锐甚。梁孝王恐。遣六将军击吴。又败梁两将。士卒皆还走梁。梁数使使报条侯求救。条侯不许。又使使恶条侯於上。上使人告条侯救梁。复守便宜不行。梁使韩安国及楚死事相弟张羽为将军。乃得颇败吴兵。吴兵欲西。梁城守坚。不敢西。即走条侯军。会下邑。欲战。条侯壁。不肯战。吴粮绝。卒饥。数挑战。遂夜饹条侯壁。惊东南。条侯使备西北。果从西北入。吴大败。士卒多饥死。乃畔散。於是吴王乃与其麾下壮士数千人夜亡去。度江走丹徒。保东越。东越兵可万馀人。乃使人收聚亡卒。汉使人以利啗东越。东越即绐吴王。吴王出劳军。即使人鏦杀吴王。盛其头。驰传以闻。吴王子子华。子驹亡走闽越。吴王之弃其军亡也。军遂溃。往往稍降太尉。梁军。楚王戊军败。自杀。\\
\hspace*{2em}三王之围齐临菑也。三月不能下。汉兵至。胶西。胶东。菑川王各引兵归。胶西王乃袒跣。席藁。饮水。谢太后。王太子德曰。﹁汉兵远。臣观之已罢。可袭。原收大王馀兵击之。击之不胜。乃逃入海。未晚也。﹂王曰。﹁吾士卒皆已坏。不可发用。﹂弗听。汉将弓高侯穨当遗王书曰。﹁奉诏诛不义。降者赦其罪。复故。不降者灭之。王何处。须以从事。﹂王肉袒叩头汉军壁。谒曰。﹁臣卬奉法不谨。惊骇百姓。乃苦将军远道至于穷国。敢请菹醢之罪。﹂弓高侯执金鼓见之。曰。﹁王苦军事。原闻王发兵状。﹂王顿首膝行对曰。﹁今者。晁错天子用事臣。变更高皇帝法令。侵夺诸侯地。卬等以为不义。恐其败乱天下。七国发兵。且以诛错。今闻错已诛。卬等谨以罢兵归。﹂将军曰。﹁王苟以错不善。何不以闻。未有诏虎符。擅发兵击义国。以此观之。意非欲诛错也。﹂乃出诏书为王读之。读之讫。曰。﹁王其自图。﹂王曰。﹁如卬等死有馀罪。﹂遂自杀。太后。太子皆死。胶东。菑川。济南王皆死。国除。纳于汉。郦将军围赵十月而下之。赵王自杀。济北王以劫故。得不诛。徙王菑川。\\
\hspace*{2em}初。吴王首反。并将楚兵。连齐赵。正月起兵。三月皆破。独赵后下。复置元王少子平陆侯礼为楚王。续元王后。徙汝南王非王吴故地。为江都王。\\
\hspace*{2em}太史公曰。吴王之王。由父省也。能薄赋敛。使其众。以擅山海利。逆乱之萌。自其子兴。争技发难。卒亡其本。亲越谋宗。竟以夷陨。晁错为国远虑。祸反近身。袁盎权说。初宠后辱。故古者诸侯地不过百里。山海不以封。﹁毋亲夷狄。以疏其属﹂。盖谓吴邪。﹁毋为权首。反受其咎﹂。岂盎。错邪。\\

% ===== 索隐述赞（司马贞《史记索隐》）· 统一样式 =====
% 题头「索隐述赞曰」五字并入反白框；正文每四句一列，定宽盒 \hfil 居列中；
% \fontsize{24}{41} 列距对齐全书界行（121+41k）；\clearpage 使述赞独占一页。
\clearpage
\shuzanhead{索隐述赞曰}\linebreak
{\shuzanbody\fontsize{24}{41}\selectfont
\hbox to \hsize{\hfil 吴楚轻悍。王濞倍德。富因采山。衅成提局。\hfil}\linebreak
\hbox to \hsize{\hfil 憍矜贰志。连结七国。婴命始监。错诛未塞。\hfil}\linebreak
\hbox to \hsize{\hfil 天之悔祸。卒取奔北。\hfil}%
\par}
\end{document}
